% Extensión estimada: unas 3.5 páginas (doce párrafos, dos figuras y una lista de tres niveles).
% Fuentes: Codigo/scripts/generate_report.py (plantillas y elección de la plantilla); Codigo/scripts/dermlip_split_eval.py
% (kNN e índice por lesión, líneas 100–111); Codigo/scripts/build_uncertainty_3tier.py (regla de ruteo);
% Codigo/scripts/ensemble_bootstrap.py (línea base sin generador); Codigo/docs/resul/revision_experimentacion_20261003.md
% (ancla Swin, índice por lesión, ruteo); Codigo/docs/ideas/rag_les_plan_20261003.md (decodificación de la recuperación);
% Codigo/docs/resul/ensamble_resul_20260607.md (piloto de noreasoning); Capítulo 2, Sección 2.5 (concordancia, consenso y
% política de tres salidas) y Sección 2.4 (kNN, Ecuación de vecinos); Figuras de Docs/experimentacion (sin cifras).
\section{Pipeline propuesto}
\label{sec:met_pipeline}

Los generadores de la Sección~\ref{sec:met_generacion} son \glspl{vlm} generativos y presentan dos limitaciones que este
trabajo evalúa por separado. Cuando clasifican, no discriminan bien las lesiones; sus resultados están en el
Capítulo~\ref{cp:res_clasif}. Cuando redactan, suelen no comprometerse con un diagnóstico y recurrir a
reservas (\emph{hedging}; Sección~\ref{sec:marco_rep_fallas}). Para atender ambas, el \emph{pipeline} separa las tareas: un
clasificador decide el diagnóstico y el generador redacta el reporte. Así, el \emph{pipeline} atiende la clasificación (T1) y la
generación del reporte (T3). La detección de melanoma (T2) queda fuera, porque el ruteo no emite una decisión binaria.

Una forma de separar las tareas es la cascada clasificador--generador. El clasificador entrega una clase, que se
inyecta en el \emph{prompt} junto con la imagen, y el generador redacta un reporte que la afirma. A esa inyección se
la llama \emph{grounding} (Sección~\ref{sec:marco_rep_generacion}). Su costo es que el reporte hereda el error del
clasificador, y su límite es que una sola clase no indica cuándo el clasificador duda (Figura~\ref{fig:met_pipeline}).

% Figura: pipeline propuesto (cascada con ruteo). Adaptada de fig:cascada y de fig:incertidumbre de
% Docs/experimentacion/experimentacion.tex, sin cifras. Se usa en sec:met_pipeline, después del párrafo que la cita.

\begin{figure}[H]
    \centering
    \begin{tikzpicture}[x=0.93mm, y=1mm]

        \node[figentrada, text width=20mm] (img) at (0, 0) {Imagen\\{\footnotesize dermatoscópica}};

        \node[figtarea, text width=32mm] (clf) at (36, 16)
            {Clasificador\\{\footnotesize Swin-Tiny}};
        \node[figtarea, text width=32mm, draw=figceleste!70!black, fill=figceleste!18] (knn) at (36, -16)
            {kNN sobre\\{\footnotesize \emph{embeddings} de DermLIP}};

        \node[fignodo, text width=24mm, fill=figrelleno] (rut) at (76, 0) {Ruteo por\\concordancia};

        \node[figtarea, text width=22mm] (gen) at (110, 0) {Generador\\{\footnotesize \emph{noreasoning}}};
        \node[figsalida, text width=14mm] (rep) at (136, 0) {Reporte};

        \node[figsalida, text width=30mm, draw=figbermellon!80!black, fill=figbermellon!12] (der) at (76, -28)
            {Derivación a\\especialista};

        \draw[figflecha] (img.east) -- ++(4, 0) |- (clf.west);
        \draw[figflecha] (img.east) -- ++(4, 0) |- (knn.west);
        \draw[figflecha] (clf.east) -| (rut.north);
        \draw[figflecha] (knn.east) -- ++(4, 0) |- (rut.west);
        \draw[figflecha] (rut.east) -- (gen.west);
        \draw[figflecha] (gen.east) -- (rep.west);
        \draw[figflecha] (rut.south) -- (der.north);

        \node[figetiqueta, anchor=west] at (78, -14) {sin consenso};

    \end{tikzpicture}
    \caption{Pipeline propuesto. Según la concordancia entre clasificador y kNN, el generador redacta un reporte
    o el caso se deriva.}
    \label{fig:met_pipeline}
\end{figure}


Un segundo predictor, con otros datos y otro sesgo, permite saber si el primero duda. En este trabajo es un \gls{knn}
sobre los \emph{embeddings} de DermLIP, que recupera los casos de entrenamiento más parecidos a la imagen
\cite{cover1967nearest}. Sus vecinos cumplen dos funciones: son evidencia textual para el generador y son la base de
una señal de concordancia con el clasificador (Sección~\ref{sec:marco_concordancia}).

\subsection{Cascada clasificador--generador (grounding)}
\label{sec:met_cascada}

La cascada depende de la instrucción con que el generador recibe el diagnóstico. La variante \emph{noreasoning} pide
solo el reporte final, prohíbe los bloques de razonamiento y prohíbe mencionar diagnósticos alternativos. El motivo es
que los generadores suelen razonar en voz alta y contradecir el diagnóstico que reciben, y prohibirles el razonamiento
de manera explícita los obliga a comprometerse con él. El \emph{grounding} en general se describe en la
Sección~\ref{sec:marco_rep_generacion}, y la plantilla completa está en el Apéndice~\ref{ap:prompts}.
La variante se eligió en un piloto con las primeras 100 imágenes de \emph{HAM10000} por orden de nombre, sin
distinguir particiones: 67 son de entrenamiento, 10 de validación, 17 de prueba y 6 quedaron fuera de la partición por
lesión. La elección se basó en la consistencia de los reportes con el diagnóstico verdadero inyectado, no en la
exactitud de ningún clasificador, y las 17 imágenes de prueba son el 1.13\,\% del test.

El diagnóstico que recibe el generador puede venir de cuatro fuentes, comparadas con el mismo generador. La primera es
Swin-Tiny, la predicción \emph{argmax} del modelo con semilla 42, que se usa como ancla principal y se entrena con la
partición por lesión (Sección~\ref{sec:met_supervisados}). La segunda es DermLIP en modo \emph{zero-shot}, una
alternativa sin entrenamiento sobre \emph{HAM10000} (Sección~\ref{sec:met_clasif_vlm}). La tercera es el \emph{oracle},
que entrega el diagnóstico verdadero: es una cota superior y no un sistema usable, porque en la práctica la etiqueta no
se conoce de antemano. La cuarta es la ausencia de anclaje: el generador sin diagnóstico, que muestra qué aporta el anclaje.

Los generadores son los de la Sección~\ref{sec:met_generacion}, con los pesos ajustados allí. La cascada no los
reentrena, de modo que el \emph{grounding} es solo de inferencia. La decodificación es la de esa sección, y los
reportes se evalúan sin el bloque de razonamiento, con la métrica de la Sección~\ref{sec:met_protocolo}.

\subsection{Grounding por recuperación}
\label{sec:met_recuperacion}

La recuperación usa el mismo kNN de la Sección~\ref{sec:met_clasif_vlm}, pero como índice de casos y no como
clasificador. Sus vecinos se calculan con \emph{embeddings} de DermLIP normalizados y similitud coseno, con $k = 3$ y
voto mayoritario. Si los tres vecinos son de clases distintas, el voto cae en el más cercano, y ese es el diagnóstico que el kNN
registra en el nivel de abstención (Sección~\ref{sec:met_ruteo}). El índice contiene solo imágenes de entrenamiento de la
partición por lesión (Sección~\ref{sec:met_particiones}), es decir, ni imágenes ni lesiones de test, y tiene 6\,447
imágenes. Esta restricción importa porque si una foto de test tiene como vecina otra de su misma lesión, el voto mide
memoria y no generalización (la fuga de información se define en la Sección~\ref{sec:marco_regularizacion}).

El contexto que recibe el generador es una frase fija con los diagnósticos de los tres vecinos y, si el metadato lo
trae, su localización. Hay dos plantillas. La de recuperación sola no usa clasificador y pide al generador el
diagnóstico. La de Swin más recuperación incorpora el diagnóstico del ancla y el contexto, sin diferenciales, y
reemplaza la plantilla \emph{noreasoning}. Por esa razón es un brazo de comparación y no entra en la cascada del ruteo.
Las corridas con recuperación generan como máximo 384 tokens, con semilla 42. En el ruteo, en cambio, la recuperación
cumple además la función de segundo predictor (Sección~\ref{sec:met_ruteo}).

\subsection{Ruteo por concordancia en tres niveles}
\label{sec:met_ruteo}

El ruteo usa la concordancia entre los dos predictores como señal. Toma sus decisiones y no sus probabilidades, así que
no depende de la calibración de la softmax ni necesita un umbral que ajustar (Sección~\ref{sec:marco_concordancia}). Los
predictores tienen sesgos distintos: uno ajusta fronteras de decisión sobre \emph{HAM10000}~\cite{tschandl2018ham10000}
y el otro no es paramétrico sobre una representación preentrenada. Por eso se espera que sus errores no coincidan
siempre. Si la concordancia separa los aciertos de los errores es una condición que se comprueba en el
Capítulo~\ref{cp:res_clasif}, no un supuesto. La concordancia se eligió en julio de 2026 con cifras de test a la vista.
La regla no tiene parámetros que ajustar, así que el test no fijó ningún umbral, pero la elección de la señal queda
declarada y los resultados del ruteo se leen con ese antecedente.

Sean $a(x)$ la concordancia entre el clasificador y el voto del kNN, y $\kappa(x)$ el consenso de los vecinos
(Ecuación~\ref{eq:marco_knn}). Con $k = 3$, el consenso vale $1/3$, $2/3$ o $1$. La regla que asigna cada caso a un
nivel es
\begin{equation}
    \pi(x) =
    \begin{cases}
        \text{afirmativo} & \text{si } a(x) = 1,\\
        \text{diferencial} & \text{si } a(x) = 0 \text{ y } \kappa(x) > 1/3,\\
        \text{abstención} & \text{si } a(x) = 0 \text{ y } \kappa(x) = 1/3.
    \end{cases}
    \label{eq:met_politica_ruteo}
\end{equation}
Los tres niveles son alternativas excluyentes, con su acción:
\begin{enumerate}
    \item \emph{Afirmativo} ($a = 1$): reporte definitivo con el diagnóstico del clasificador y la plantilla
    \emph{noreasoning}.
    \item \emph{Diferencial} ($a = 0$ y $\kappa > 1/3$): reporte que enuncia el par formado por el clasificador y el
    voto del kNN, con la plantilla diferencial.
    \item \emph{Abstención} ($a = 0$ y $\kappa = 1/3$, es decir, los tres vecinos de clases distintas): no se genera
    reporte y el caso se deriva a un especialista.
\end{enumerate}
Las plantillas están en el Apéndice~\ref{ap:prompts}. La plantilla diferencial atribuye la incertidumbre al
clasificador, aunque uno de los candidatos venga del kNN, y se describe tal cual está. Como la señal es discreta, no hay
umbral que ajustar, y los tres niveles cubren todos los casos (Figura~\ref{fig:met_ruteo}).

% Figura: ruteo en tres niveles según la concordancia. Adaptada de fig:incertidumbre de
% Docs/experimentacion/experimentacion.tex; se quitan los porcentajes y las coberturas. Se usa en sec:met_ruteo.

\begin{figure}[H]
    \centering
    \begin{tikzpicture}[x=1mm, y=1mm]

        \node[figentrada, text width=30mm] (sig) at (0, 0)
            {Clasificador y kNN\\{\footnotesize concordancia $a(x)$}};

        \node[figsalida, text width=36mm, draw=figverde!70!black, fill=figverde!12] (aff) at (80, 24)
            {Afirmativo\\{\footnotesize reporte definitivo}};
        \node[figsalida, text width=36mm, draw=fignaranja!70!black, fill=fignaranja!18] (dif) at (80, 0)
            {Diferencial\\{\footnotesize reporte con dos candidatos}};
        \node[figsalida, text width=36mm, draw=figbermellon!80!black, fill=figbermellon!12] (abs) at (80, -24)
            {Derivación\\{\footnotesize sin reporte, a especialista}};

        \draw[figflecha] (sig.east) -- (30, 0) |- (aff.west);
        \draw[figflecha] (30, 0) |- (abs.west);
        \draw[figflecha] (sig.east) -- (dif.west);

        \node[figetiqueta, above] at (45, 24) {$a = 1$};
        \node[figetiqueta, above] at (45, 0) {$a = 0$ y $\kappa > 1/3$};
        \node[figetiqueta, below] at (45, -24) {$a = 0$ y $\kappa = 1/3$};

    \end{tikzpicture}
    \caption{Ruteo en tres niveles según la concordancia entre el clasificador y el kNN.}
    \label{fig:met_ruteo}
\end{figure}


Esta regla es la política de tres salidas de la Sección~\ref{sec:marco_tres_salidas}, con la concordancia en lugar de la
confianza. La abstención es una predicción selectiva (Sección~\ref{sec:marco_pred_selectiva}): en ese nivel el sistema
no emite reporte. El diferencial presenta como candidatos el diagnóstico del clasificador y el voto del kNN, y no los
dos primeros del clasificador, porque lo que el ruteo mide es el desacuerdo entre ambos predictores. La regla dice
qué se emite, pero no cuánto aporta el generador.

La línea base sin generador aísla ese aporte. Usa el mismo ruteo y el mismo clasificador, pero no redacta: el nivel
afirmativo devuelve el diagnóstico del clasificador, el diferencial devuelve el par de candidatos y el de abstención
queda igual que en el sistema. La cascada se compara con líneas base que usan el mismo clasificador, y la
Sección~\ref{sec:met_complementarios} describe las pruebas con que se analiza esa diferencia.
